\chapter{Development of Concepts}
\label{ch:development-concepts}

\section{Purpose and scope of concept development}
Concept development is an essential part of structured design, which consists of the design phase where different principles solutions are developed and evaluated before implementation. This is an essential part of systematic engineering design. By evaluating different options with respect to certain criteria before proceeding, the chances of selecting an inappropriate design path are minimized. REF The concept development phase is crucial for the success of the project, as it allows for the exploration of different design alternatives and the selection of the most suitable one based on a structured evaluation process.
\par The gap in the research, as identified in Section \ref{Summary of the State of the Art and Identified Research Gap}, reveals the absence of an effective systematic approach to automatically generate robot programs for press brake sheet metal bending. This approach should consider changes in workpiece geometry, gripping, re-gripping, and pre-execution pose, minimizing manual programming. The gap in the research is addressed in the form of measurable engineering objectives, as discussed in \ref{CH_3_System Concept and Methodology}, with respect to functional requirements FR-01 to FR-07 and non-functional requirements NFR-01 to NFR-05. This chapter presents three different strategies to fulfill these objectives.
\par Concept A proposes a guided teaching approach in which a structured software workflow supports the operator during conventional manual robot programming. Concept B proposes an automatic program generation system in which waypoint, grip positions , regrip sequences and opposite side handling strategies are derived automatically by evaluating the inputs configuration file , with communication to the robot runtime via WebSocket.\Cite{dragandbot_anyrobot_2025} Concept C is a CAD integrated offline programming in which a digital cell model is used to generate , validate and execute robot program within a simulation environment \cite{Siciliano2016Robotics},\cite{Pan2012Survey}
\par The concept decision is made using Weighted Sum Model (WSM), a quantitative decision support method in which evaluation criteria are assigned relative weights and each concept is scored against criterion. This method satisfies the requirement that concept selection must be based on structure rather than on qualitative judgment. The morphological box and preferences analysis methods are alternative methods, The WSM is selected because it enables simultaneous multi-criteria weighting and produces a single quantitative score for each concept, which ensures appropriate for comparing concepts.
%%%%%%%%%%%%%%%%%%%%%%%%%%%%%%
\section{Concept A – Structured operator Workflow with Guided Teach in programming}
\label{Concept A – Structured operator Workflow with Guided Teach in programming}
\subsection{Methodological Framework}
Concept A retains the conventional teach pendant paradigm as the mechanism for defining robot motion and introduces a software-based layer that structures and sequences the operator interaction. In this method, the operator is guided through a step-by-step workflow that collects process inputs, including the desires bending sequences, bend side selection and current workpiece positions. The operator then executes each action directly via robot teach pendant or control interface using the function blocks.
\par This Concept is grounded in the established practice of teach pendant programming described in section\ref{Manual And teach pendant based robot programming}. As identified in the literature, teach pendant programming is reliable and widely used in structured manufacturing environment where part geometries and process parameters are fixed \cite{Siciliano2016Robotics}. The guidance layer introduced in this concept (A) reduced the risk of operator errors and provides a structured interface than unguided teach pendant operation, but it does not alter the fundamental that all robot position and waypoints are defined through direct physical interaction with the robot.
\par Regripping operation, orientation changes, full side bending operation for a rectangular sheet under this concept A remains the responsibility of the user to define and execute. The system provides instructional prompts but does not execute, validate or generate any waypoints and position data. No interface to the robot runtime environment is required beyond standard teach pendant connectivity. The system does not possess object tracking subsystems, evaluate geometric constraint, grip position calculations or assess kinematic feasibility before execution.

\subsection{System Structure}
\label{System Structure}
The system structure of concept A consists of two layers. A dialogue management layer runs on an external computer and presents the operator prompt through functional blocks or graphical interface by manual program generation. A robot controller layer receives waypoint generated manually by the operator through teach pendant and execute them. No intermediate computation layer for geometric reasoning, constraint evaluation, or motion generation exists between these layers. The architecture imposes no dependency on object tracking sensor data or trajectory generation interfaces.

\subsection{Critical Evaluation}
\label{Critical Evaluation}
Concept A straightforward to implement and compatible with any industrial robot controller without software modification. The operation of the system is fully transparent and observable by the operator, which reduces the risk of unreachable and kinematics error occur during the unreachable target pose. These characteristics make it suitable for low variety, structured production scenario where process parameters are stable across batches.
\par However , the fundamentals limitation if this concept are significant in the context of research gap ,programming effort is not yet reduced because knowledge of robot kinematics , tooling constraints , regripping strategies and bending sequence algorithm should need to applied by the operator for every session \cite{Siciliano2016Robotics}, \cite{Pan2012Survey}. Minor changes in sheet geometry , workpiece holder placement requires full re teaching of the affected trajectories , which contradicts the flexible objective stated  in section \ref{Problem-Definition-and-Objectives} \cite{Siciliano2016Robotics}, REF 26. Regripping operation and full side bending operations must ne manually planned , which introduces a chances of operator error and increase the cycle time .No pre execution validation of target poses is performed , kinematic configuration are  discovered only during the physical or virtual robot motion which leads to safety risk and a source of unplanned downtime \cite{Siciliano2016Robotics}.
\par The software guidance layer of concept A provides process management but no automation level, it addresses the unstructured operator workflow without addressing root cause, which is the requirement for manual position definition. Concept A does not contribute to the primary objective of automatic program generation and failed to satisfy most of the functional and non-function requirements defined in \ref{tab:functional_requirements}. and REF NFR.

%%%%%%%%%%%%%%%%%%%%
\section{Concept B – Automatic Program Generation based on rules with Integrated Constraint Validation.}
\label{Concept B – Automatic Program Generation based on rules with Integrated Constraint Validation.}
\subsection{Methodological Framework}
Concept B proposes an automation level, rule based, automated program for generation system in which all robot waypoints, grip positions, regrip sequences, opposite side handling operations and strategies are derives automatically from geometric input parameters and a structures YAML configuration file: operator interaction is limited selection of required output selection before execution and confirmation of generated programs visually. No manual waypoint generation, direct teaching or any runtime operation or decisions are required during program execution.
\par The system communicates with the robot runtime through AnyRobot interface, a robotics application program (API) Developed by KEBA \textbackslash\ Drag and Bot GmbH that provides standardized method for controlling industrial robots from different manufacturers via WebSocket communication\cite{dragandbot_anyrobot_2025}. The planning Software operates as the WebSocket client, the drag and bot runtime act as the server. This architecture establishes a clean separation between the planning computation layer and robot execution layer, which is a recognized design principle \cite{Thramboulidis2013MDE} .
\par Workpiece pose data is retrieved in real time from drag and bot object tracking subsystem and processes by an orientation validation component , which determines whether the tracked sheet orientation is admissible for the planned bending operation by applying a three level priority strategy ,a side normalized orientation is evaluation first , followed by original tracked orientations by the object tracker and finally a 180 degree rotated strategy , when neither of the first two falls within the configurable window. This Three level hierarchy ensures that valid gripping orientations and positions are identified for any initial sheet placement on the cell floor, which eliminates the need for operator prepositioning. Before grip position calculation, the program validated the sheet position is within the robot cell. This is applied to avoid the kinematic reachability error occurred during the program execution. The zone-specific admissibility limits are specific limits are robot cells configuration are configured in the input sections and require no code modification when cell layout changes. This design decision satisfies NFR 02as defined in table \ref{tab:functional_requirements} and NFR 05 as defined in table REF TABLE NFR .
\par Grip position are computed by a parameterized grip planning component s that uses the sheet dimension, mass suction cup diameter , rated shear force and edge clearance margin which all are defined in the input configuration files :candidate grip position reevaluated against edge clearance and cup spacing constraints and against the vacuum holding force which is derived from the relationship between suction pressure , cup contact areas and dynamic loading during transport \cite{VacuumGripperModeling}.Grip feasibility is evaluated before kinematic reachability assessment, ensuring that only variable grip configurations are submitted to the motion generation layer. This ordering prevents unsafe handling states that would result from a grip position that is kinematically reachable but physically unable to reach the sheet during transport \cite{VacuumGripperModeling} .
\par When the current grip configuration does not allow the targeted pose to be reached, an automated regrip cycle is initiated. The cycle places the sheet on a designated orientation table, repositions the robot to reset and then regrasp the sheet at an orientation that restores full reachability. This regrips logic directly satisfies FR03 as defined in table \ref{tab:functional_requirements} .
\par Opposite side bending is also supported through this concept through a dedicated module that executes a structured sheet flip sequence. After all primary bending operations are completed, the program evaluates whether the operated has selected for opposite side bends. If so, the sheet passes to flip position through a coordinated sequence of motion, gripper activation, and object attachment operations targeting positions to flip the sheet. Afte the flip, the opposite bending sequence is executed under the same validation, including orientation normalization, grip feasibility evaluation and reachability assessment before the waypoints is generated. This enforcement of validation across both primary and opposite side bending operation ensures consistent safety behavior throughout th full step bending process \cite{Gupta1998SheetMetalBending}
\par The motion trajectories are generated in linear cartesian motion commands using minimum jerk velocity interpolation profile REF 16: This profile minimizes the third derivatives of position over the motion, producing smooth velocity transition between waypoint that reduce peak acceleration and gripper loading forces during sheet transport. This is requirement for maintaining vacuum grip stability under dynamic conditions \cite{VacuumGripperModeling}.

%%%%%%%%%%%%%%%%%%%%%%%%
\section{Concept C – CAD Integrated Offline Programming with Digital Twin Simulation}
\label{Concept C – CAD Integrated Offline Programming with Digital Twin Simulation}
\subsection{Methodological Framework}